
\documentclass[11pt,a4paper]{ctexart}
\usepackage[margin=2.2cm]{geometry}
\usepackage{xcolor}
\usepackage{listings}
\usepackage{booktabs}
\usepackage{longtable}
\usepackage{enumitem}
\usepackage{hyperref}
\usepackage{fancyhdr}

\hypersetup{colorlinks=true,linkcolor=blue!60!black,urlcolor=blue!60!black}
\pagestyle{fancy}
\fancyhf{}
\fancyhead[L]{Statistical Programming II}
\fancyhead[R]{Practical 6: R Packages}
\fancyfoot[C]{\thepage}

\definecolor{codebg}{RGB}{247,247,247}
\definecolor{codeframe}{RGB}{210,210,210}
\lstdefinelanguage{R}{
  morekeywords={function,if,else,TRUE,FALSE,NULL,NA},
  otherkeywords={|>,::,<-,\%in\%},
  sensitive=true,
  morecomment=[l]{\#},
  morestring=[b]",
  morestring=[b]'
}
\lstset{
  language=R,
  backgroundcolor=\color{codebg},
  frame=single,
  rulecolor=\color{codeframe},
  basicstyle=\ttfamily\small,
  breaklines=true,
  showstringspaces=false,
  columns=fullflexible,
  keepspaces=true
}

\title{\textbf{Statistical Programming II}\\Practical 6: R Packages\\题目与答案整理}
\author{Youpeng Jiang}
\date{\today}

\begin{document}
\maketitle
\tableofcontents
\newpage

\section{总体目标}

本次 Practical 从零创建 R Package \texttt{munichvisitors}，包括：

\begin{itemize}
  \item Package 基础结构、\texttt{DESCRIPTION}、README 和许可证；
  \item 下载、清理并保存慕尼黑博物馆访客数据；
  \item 编写、记录并导出 \texttt{plot\_museums()}；
  \item 使用 testthat 测试并用 \texttt{devtools::check()} 检查；
  \item 为第二个慕尼黑开放数据集设计带参数的新函数。
\end{itemize}

\section{Part 1：创建 Package}

\subsection{Exercise 1.1：初始化}

在 RStudio Console 中运行：

\begin{lstlisting}
usethis::create_package("munichvisitors")
\end{lstlisting}

也可以使用绝对路径：

\begin{lstlisting}
usethis::create_package(
  "C:/Data/StatProg_practical/practical_6/munichvisitors"
)
\end{lstlisting}

创建后的核心结构：

\begin{verbatim}
munichvisitors/
├── DESCRIPTION
├── NAMESPACE
├── R/
├── munichvisitors.Rproj
├── .gitignore
└── .Rbuildignore
\end{verbatim}

\begin{itemize}
  \item \texttt{DESCRIPTION}：Package 元数据和依赖。
  \item \texttt{NAMESPACE}：控制导出和导入，通常由 roxygen2 生成。
  \item \texttt{R/}：存放 Package 正式函数。
  \item \texttt{.gitignore}：Git 忽略规则。
  \item \texttt{.Rbuildignore}：构建 Package 时忽略的文件。
\end{itemize}

\subsection{Exercise 1.2：DESCRIPTION}

在项目根目录的 \texttt{DESCRIPTION} 中写：

\begin{verbatim}
Package: munichvisitors
Title: Monthly Visitor Counts for Munich Museums
Version: 0.0.0.9000
Authors@R:
    person("Youpeng", "Jiang", , "your-email@example.com",
           role = c("aut", "cre"))
Description: Provides tidy monthly visitor statistics for Munich's museums
    from Munich Open Data, with a helper plot function.
License: MIT + file LICENSE
Encoding: UTF-8
Roxygen: list(markdown = TRUE)
RoxygenNote: 7.3.2
LazyData: true
Imports:
    dplyr,
    ggplot2,
    scales
Suggests:
    janitor,
    readr,
    testthat (>= 3.0.0),
    knitr,
    rmarkdown
Config/testthat/edition: 3
VignetteBuilder: knitr
\end{verbatim}

延续行必须缩进。\texttt{Imports} 是运行时依赖，\texttt{Suggests} 是测试、文档或数据预处理依赖。

添加 MIT License：

\begin{lstlisting}
usethis::use_mit_license()
\end{lstlisting}

\subsection{Exercise 1.3：README}

Console：

\begin{lstlisting}
usethis::use_readme_rmd()
\end{lstlisting}

\texttt{README.Rmd} 的基本结构：

\begin{verbatim}
---
output: github_document
---

# munichvisitors

The package provides tidy monthly visitor counts for Munich museums.

## Installation

```r
devtools::install_github("your-username/munichvisitors")
```

## Included objects

- `museum_visitors`
- `plot_museums()`

## Example

```r
library(munichvisitors)
plot_museums()
```
\end{verbatim}

渲染：

\begin{lstlisting}
rmarkdown::render("README.Rmd")
\end{lstlisting}

主要编辑 \texttt{README.Rmd}，生成的 \texttt{README.md} 供 GitHub 显示。

\section{Part 2：添加数据}

\subsection{Exercise 2.1：data-raw}

Console：

\begin{lstlisting}
usethis::use_data_raw("museum-visitors")
\end{lstlisting}

生成：

\begin{verbatim}
data-raw/museum-visitors.R
\end{verbatim}

\texttt{data-raw/} 保存生成数据的脚本，\texttt{data/} 保存最终数据对象。

\subsection{Exercise 2.2：下载、清理和保存}

Console 中声明开发依赖：

\begin{lstlisting}
usethis::use_package("readr", type = "Suggests")
usethis::use_package("janitor", type = "Suggests")
\end{lstlisting}

在 \texttt{data-raw/museum-visitors.R} 中写：

\begin{lstlisting}
url <- paste0(
  "https://opendata.muenchen.de/dataset/",
  "bfb4a286-bea5-4bfe-82ce-b9bd354284a5/resource/",
  "6c6a809e-91ee-4f3e-9268-a8b7bc38311c/download/",
  "monatszahlen2603_museen_16_03_26.csv"
)

museum_visitors_raw <- readr::read_csv(url)

museum_visitors <- museum_visitors_raw |>
  janitor::clean_names()

usethis::use_data(
  museum_visitors,
  overwrite = TRUE
)
\end{lstlisting}

执行整个脚本后生成：

\begin{verbatim}
data/museum_visitors.rda
\end{verbatim}

\subsection{Exercise 2.3：检查数据}

Console：

\begin{lstlisting}
devtools::load_all()

museum_visitors
dplyr::glimpse(museum_visitors)
names(museum_visitors)
dim(museum_visitors)
class(museum_visitors)
\end{lstlisting}

重要列包括：

\begin{verbatim}
monatszahl
auspraegung
jahr
monat
wert
vorjahreswert
veraend_vormonat_prozent
veraend_vorjahresmonat_prozent
zwoelf_monate_mittelwert
\end{verbatim}

自动检查：

\begin{lstlisting}
expected_columns <- c(
  "monatszahl", "auspraegung", "jahr", "monat", "wert",
  "vorjahreswert", "veraend_vormonat_prozent",
  "veraend_vorjahresmonat_prozent",
  "zwoelf_monate_mittelwert"
)

all(expected_columns %in% names(museum_visitors))
\end{lstlisting}

\subsection{Exercise 2.4：数据文档}

Console：

\begin{lstlisting}
usethis::use_r("data")
\end{lstlisting}

在 \texttt{R/data.R} 中写：

\begin{lstlisting}
#' Monthly visitor counts for Munich museums
#'
#' Monthly and annual visitor statistics for Munich's public museums.
#'
#' @format A data frame with monthly and annual observations:
#' \describe{
#'   \item{monatszahl}{Category label}
#'   \item{auspraegung}{Museum name}
#'   \item{jahr}{Year}
#'   \item{monat}{YYYYMM code or "Summe"}
#'   \item{wert}{Monthly value or annual total}
#'   \item{vorjahreswert}{Value for the prior period}
#'   \item{veraend_vormonat_prozent}{Change vs. previous month}
#'   \item{veraend_vorjahresmonat_prozent}{Change vs. prior year}
#'   \item{zwoelf_monate_mittelwert}{12-month rolling average}
#' }
#' @source Landeshauptstadt München. Monatszahlen Museen.
"museum_visitors"
\end{lstlisting}

生成并查看帮助：

\begin{lstlisting}
devtools::document()
devtools::load_all()
?museum_visitors
\end{lstlisting}

数据文档最后使用字符串 \texttt{"museum\_visitors"}，不需要 \texttt{@export}。

\section{Part 3：添加函数}

\subsection{Exercise 3.1：plot\_museums()}

Console：

\begin{lstlisting}
usethis::use_r("plot_museums")
usethis::use_package("dplyr")
usethis::use_package("ggplot2")
usethis::use_package("scales")
\end{lstlisting}

在 \texttt{R/plot\_museums.R} 中写：

\begin{lstlisting}
#' Annual visitor counts per Munich museum
#'
#' Plots annual visitor totals using rows where `monat == "Summe"`.
#'
#' @return A `ggplot2` plot object.
#' @export
#' @examples
#' plot_museums()
plot_museums <- function() {
  museum_visitors |>
    dplyr::filter(monat == "Summe") |>
    ggplot2::ggplot(
      ggplot2::aes(
        x = jahr,
        y = wert,
        colour = auspraegung
      )
    ) +
    ggplot2::geom_line() +
    ggplot2::labs(
      title = "Annual Visitors to Museums in Munich",
      x = "Year",
      y = "Visitors",
      colour = "Museum"
    ) +
    ggplot2::scale_y_continuous(
      labels = scales::label_number()
    )
}
\end{lstlisting}

虽然题目有时称其为 monthly plot，但筛选 \texttt{monat == "Summe"} 后实际画的是年度总数。

创建 \texttt{R/utils.R}：

\begin{lstlisting}
usethis::use_r("utils")
\end{lstlisting}

内容：

\begin{lstlisting}
utils::globalVariables(
  c(
    "museum_visitors",
    "monat",
    "wert",
    "auspraegung",
    "jahr"
  )
)
\end{lstlisting}

\subsection{Exercise 3.2：文档和导出}

运行：

\begin{lstlisting}
devtools::document()
devtools::load_all()

plot_museums()
?plot_museums
\end{lstlisting}

生成：

\begin{verbatim}
man/plot_museums.Rd
NAMESPACE
\end{verbatim}

\texttt{NAMESPACE} 中应有：

\begin{verbatim}
export(plot_museums)
\end{verbatim}

这来自 \texttt{\#' @export}，不要手动修改 \texttt{NAMESPACE}。

\subsubsection*{Usage 来自哪里？}

帮助页面中的：

\begin{verbatim}
Usage

plot_museums()
\end{verbatim}

由函数定义自动生成：

\begin{lstlisting}
plot_museums <- function() {
\end{lstlisting}

如果定义是：

\begin{lstlisting}
plot_museums <- function(museum, start_year = 2020) {
\end{lstlisting}

Usage 会自动显示：

\begin{verbatim}
plot_museums(museum, start_year = 2020)
\end{verbatim}

\subsection{Exercise 3.3：Package check}

Console：

\begin{lstlisting}
devtools::check()
\end{lstlisting}

目标：

\begin{verbatim}
0 errors | 0 warnings | 0 notes
\end{verbatim}

常见问题：

\begin{itemize}
  \item 缺少 \texttt{@export}；
  \item 外部函数未写成 \texttt{package::function()}；
  \item 依赖未写入 \texttt{DESCRIPTION}；
  \item 裸列名未在 \texttt{globalVariables()} 中声明；
  \item 函数或数据缺少文档；
  \item Examples 在干净环境中不能运行。
\end{itemize}

\subsection{Exercise 3.4：自动测试}

Console：

\begin{lstlisting}
usethis::use_testthat()
usethis::use_test("plot_museums")
\end{lstlisting}

在 \texttt{tests/testthat/test-plot\_museums.R} 中写：

\begin{lstlisting}
test_that("plot_museums returns a ggplot object", {
  result <- plot_museums()
  expect_s3_class(result, "gg")
})

test_that("museum_visitors has expected columns", {
  expect_true(
    all(
      c("auspraegung", "jahr", "monat", "wert") %in%
        names(museum_visitors)
    )
  )
})
\end{lstlisting}

运行：

\begin{lstlisting}
devtools::test()
devtools::check()
\end{lstlisting}

\section{Part 4：第二个数据集}

这里选择慕尼黑 Theater 数据，设计：

\begin{lstlisting}
get_theatre_stats(
  indicator = "Besucher*innen",
  years = NULL
)
\end{lstlisting}

\subsection{创建文件和依赖}

Console：

\begin{lstlisting}
usethis::use_r("get_theatre_stats")
usethis::use_package("readr", type = "Imports")
usethis::use_package("janitor", type = "Imports")
\end{lstlisting}

因为函数运行时在线读取和清理数据，\texttt{readr} 和 \texttt{janitor} 属于 \texttt{Imports}。

\subsection{函数代码}

在 \texttt{R/get\_theatre\_stats.R} 中写：

\begin{lstlisting}
#' Download monthly statistics for Munich theatres
#'
#' @param indicator One of `"Besucher*innen"`, `"Aufführungen"`,
#'   or `"Platzausnutzung"`.
#' @param years Optional numeric vector of years. `NULL` keeps all.
#'
#' @return A tibble containing monthly theatre statistics.
#' @export
#'
#' @examples
#' \dontrun{
#' get_theatre_stats()
#' get_theatre_stats("Aufführungen", 2019:2024)
#' }
get_theatre_stats <- function(
    indicator = "Besucher*innen",
    years = NULL
) {
  allowed_indicators <- c(
    "Besucher*innen",
    "Aufführungen",
    "Platzausnutzung"
  )

  if (
    !is.character(indicator) ||
      length(indicator) != 1L ||
      is.na(indicator) ||
      !indicator %in% allowed_indicators
  ) {
    stop(
      paste0(
        "`indicator` must be one of: ",
        paste(allowed_indicators, collapse = ", "),
        "."
      ),
      call. = FALSE
    )
  }

  if (
    !is.null(years) &&
      (!is.numeric(years) || anyNA(years))
  ) {
    stop(
      "`years` must be NULL or a numeric vector without missing values.",
      call. = FALSE
    )
  }

  url <- paste0(
    "https://opendata.muenchen.de/dataset/",
    "dacf1009-8933-4cdf-8c9d-40afd54e4ef7/resource/",
    "eb734a10-3b0a-4421-ae1d-4226f9409538/download/",
    "monatszahlen2510_theater_30_10_25.csv"
  )

  theatre_data <- readr::read_csv(
    url,
    show_col_types = FALSE
  ) |>
    janitor::clean_names()

  result <- theatre_data |>
    dplyr::filter(
      monatszahl == indicator,
      monat != "Summe",
      !is.na(wert)
    )

  if (!is.null(years)) {
    result <- result |>
      dplyr::filter(jahr %in% years)
  }

  result
}
\end{lstlisting}

在 \texttt{R/utils.R} 中加入 \texttt{monatszahl}：

\begin{lstlisting}
utils::globalVariables(
  c(
    "museum_visitors",
    "monatszahl",
    "monat",
    "wert",
    "auspraegung",
    "jahr"
  )
)
\end{lstlisting}

\subsection{手动检查代码放在哪里？}

以下内容直接输入 RStudio Console，不放入 \texttt{R/}：

\begin{lstlisting}
devtools::document()
devtools::load_all()

theatre_visitors <- get_theatre_stats()
dplyr::glimpse(theatre_visitors)
unique(theatre_visitors$monatszahl)

theatre_2019_2024 <- get_theatre_stats(
  indicator = "Aufführungen",
  years = 2019:2024
)

range(theatre_2019_2024$jahr)
unique(theatre_2019_2024$monatszahl)
unique(theatre_2019_2024$monat)
\end{lstlisting}

\begin{longtable}{p{0.43\textwidth}p{0.47\textwidth}}
\toprule
代码 & 放置位置 \\
\midrule
函数定义 & \texttt{R/get\_theatre\_stats.R} \\
\texttt{globalVariables()} & \texttt{R/utils.R} \\
\texttt{devtools::*()}、\texttt{usethis::*()} & RStudio Console \\
手动调用和检查结果 & RStudio Console \\
\texttt{test\_that()}、\texttt{expect\_*()} & \texttt{tests/testthat/} \\
\bottomrule
\end{longtable}

\subsection{测试第二个函数}

Console：

\begin{lstlisting}
usethis::use_test("get_theatre_stats")
\end{lstlisting}

在 \texttt{tests/testthat/test-get\_theatre\_stats.R} 中写：

\begin{lstlisting}
test_that("invalid theatre indicators are rejected", {
  expect_error(
    get_theatre_stats(indicator = "Kinobesuche"),
    "`indicator` must be one of"
  )
})

test_that("years must be numeric", {
  expect_error(
    get_theatre_stats(years = "2024"),
    "`years` must be NULL or a numeric vector"
  )
})
\end{lstlisting}

最后运行：

\begin{lstlisting}
devtools::document()
devtools::test()
devtools::check()
\end{lstlisting}

\section{常用命令速查}

\begin{longtable}{p{0.43\textwidth}p{0.47\textwidth}}
\toprule
命令 & 作用 \\
\midrule
\texttt{usethis::create\_package()} & 创建 Package。\\
\texttt{usethis::use\_r()} & 创建 \texttt{R/} 源文件。\\
\texttt{usethis::use\_data\_raw()} & 创建数据预处理脚本。\\
\texttt{usethis::use\_data()} & 保存 Package 数据。\\
\texttt{usethis::use\_package()} & 声明依赖。\\
\texttt{usethis::use\_testthat()} & 初始化测试。\\
\texttt{usethis::use\_test()} & 创建测试文件。\\
\texttt{devtools::document()} & 生成帮助文件和 NAMESPACE。\\
\texttt{devtools::load\_all()} & 加载开发版本。\\
\texttt{devtools::test()} & 运行测试。\\
\texttt{devtools::check()} & 完整检查 Package。\\
\bottomrule
\end{longtable}

\section{Deutsch: Kurzzusammenfassung}

Das Practical entwickelt ein vollständiges R-Package namens
\texttt{munichvisitors}. Konsolenbefehle wie
\texttt{usethis::use\_*()}, \texttt{devtools::document()},
\texttt{devtools::load\_all()}, \texttt{devtools::test()} und
\texttt{devtools::check()} werden in der RStudio-Konsole ausgeführt.
Funktionsdefinitionen gehören in \texttt{R/}, Datenskripte in
\texttt{data-raw/} und Tests in \texttt{tests/testthat/}.

Der vollständige Workflow lautet:

\begin{verbatim}
Code oder Dokumentation ändern
        ↓
devtools::document()
        ↓
devtools::load_all()
        ↓
devtools::test()
        ↓
devtools::check()
\end{verbatim}

Das Ziel ist ein Check mit:

\begin{verbatim}
0 errors | 0 warnings | 0 notes
\end{verbatim}

\end{document}
